\section{Matched phases and coefficient extraction}
\label{sec:app_ft_matched_phases}

This section isolates two elementary steps used by the main proof. First, the
sector partners obtained in both directions determine a finite bijection.
Second, normalization of the matched normal tensors forces every gauge scalar
to be a phase.

\begin{theorem}[Bijective matching under eventual proportionality]
    \label{thm:sector_bnt_proportional_bijective_match}
    \lean{MPSTensor.bijective_match_of_eventuallyProportional}
    \leanok
    \uses{def:sector_bnt_cf, thm:sector_bnt_proportional_full_basis_match,
        lem:sector_bnt_bijection_from_matches}
    Let $P$ and $Q$ be BNT canonical forms whose total MPV families are
    eventually nonzero and proportional. There is a bijection
    $\beta:J(Q)\simeq J(P)$ such that every $Q_k$ has an equal-dimensional,
    gauge-phase-equivalent partner $P_{\beta(k)}$, and their cross-overlap does
    not tend to zero.
\end{theorem}

\begin{proof}\leanok
    Apply the full-basis matching theorem first from $Q$ to $P$ and then to the
    reverse proportionality from $P$ to $Q$. If two sectors had the same
    partner, transitivity of gauge-phase equivalence would contradict BNT basis
    separation. Both partner maps are therefore injective. The two finite
    sector sets have equal cardinality, so the forward map is bijective and
    retains its dimension, gauge, and nondecay properties.
\end{proof}

\begin{lemma}[Unit phase from matched normalized BNT blocks]%
    \label{lem:sector_bnt_norm_phase_of_matched_mpv}
    \lean{MPSTensor.IsBNTCanonicalForm.norm_phase_of_matched_mpv}
    \leanok
    \uses{def:sector_bnt_cf}
    Let $P$ and $Q$ be BNT canonical forms. If, for every positive length $N$
    and word $\sigma$, a sector $P_j$ and a sector $Q_k$ have MPV families
    related by
    \begin{align}
        V^{(N)}(Q_k)_\sigma
         & = \zeta^N V^{(N)}(P_j)_\sigma,
        \label{eq:ft_matched_mpv_phase_norm}
    \end{align}
    then $|\zeta|=1$.
\end{lemma}

\begin{proof}\leanok
    The BNT normalization gives self-overlap limits of modulus $1$ for both
    sectors. The MPV identity (\ref{eq:ft_matched_mpv_phase_norm}) gives, for
    every positive $N$,
    \begin{align}
        O_{Q_kQ_k}(N)
         & = |\zeta|^{2N}\,O_{P_jP_j}(N).
        \notag
    \end{align}
    Since $O_{Q_kQ_k}(N)\to1$ and $O_{P_jP_j}(N)\to1$, taking limits forces
    $|\zeta|^{2N}\to1$ and hence $|\zeta|=1$.
\end{proof}


\section{Coefficient identities from gauge data}
\label{sec:app_ft_coefficient_wrappers}

The main coefficient theorem takes the matched phases as data. The next result
shows how those phases are selected from gauge-phase equivalences before the
linear-independence argument is applied.

\begin{lemma}[Full-basis coefficient identity from matched sectors]%
    \label{lem:sector_bnt_coeff_identity_via_global_gauge}
    \lean{MPSTensor.coeff_identity_via_global_gauge}
    \leanok
    \uses{def:sector_bnt_cf,
        def:gauge_phase_equiv}
    Let $P$ and $Q$ be BNT canonical forms with the same MPV family at every
    positive length, matched by a bijection $\beta:J(Q)\simeq J(P)$ with
    bond-dimension equalities and a
    gauge-phase equivalence between $Q_k$ and $P_{\beta(k)}$ for every $k$. Then
    each sector carries a unit-modulus phase $\zeta_k$, and eventually
    \begin{align}
        c_N^{(\beta(k))}(P)
         & = \zeta_k^N\,c_N^{(k)}(Q).
        \label{eq:ft_coeff_identity}
    \end{align}
    This is the equal-MPV coefficient comparison of
    \cite[Appendix MPV proof, lines~1187--1188]{Cirac2016MPDO_arXiv}.
\end{lemma}

\begin{proof}\leanok
    \uses{lem:sector_bnt_norm_phase_of_matched_mpv,
        lem:sector_bnt_coeff_identity_via_matched_mpv_phase}
    For each matched sector, the gauge-phase equivalence gives
    $V^{(N)}(Q_k)_\sigma=\zeta_k^NV^{(N)}(P_{\beta(k)})_\sigma$.
    Lemma~\ref{lem:sector_bnt_norm_phase_of_matched_mpv} gives
    $|\zeta_k|=1$. Applying
    Lemma~\ref{lem:sector_bnt_coeff_identity_via_matched_mpv_phase} to the
    same phases gives
    $c_N^{(\beta(k))}(P)=\zeta_k^Nc_N^{(k)}(Q)$ for all sufficiently
    large $N$.
\end{proof}


\section{The scalar-threaded proportional identity}
\label{sec:app_ft_proportional_scalar}

For proportional total MPVs, coefficient extraction retains the one scalar
that relates the two total vectors at each length. This identity is exact, but
it does not imply the scalar-free identity needed to recover copy weights; the
counterexample in Remark~\ref{rem:proportional_length_scalar_gap} shows that no
such implication is possible.

\begin{lemma}[Proportional matched-phase coefficient identity, scalar-threaded]%
    \label{lem:sector_bnt_coeff_identity_via_matched_mpv_phase_proportional}
    \lean{MPSTensor.coeff_identity_via_matched_mpv_phase_proportional}
    \leanok
    \uses{def:sector_bnt_cf, lem:eventual_coeff_eq_of_eventual_li}
    Let $P$ be a BNT canonical form and $Q$ a sector decomposition whose total
    tensors generate eventually nonzero proportional MPV families. Suppose a
    bijection $\beta:J(Q)\simeq J(P)$ and scalars $\zeta_k$ satisfy the matched
    MPV phase identities (\ref{eq:ft_phase_relation_hyp}) for every sector $k$.
    Then there is a single eventually-nonzero scalar sequence $c_N$ and a
    threshold $N_0$ such that, for all $N>N_0$ and every $k\in J(Q)$,
    \begin{align}
        c_N^{(\beta(k))}(P)
         & = c_N\,\zeta_k^N\,c_N^{(k)}(Q).
        \label{eq:ft_coeff_identity_threaded}
    \end{align}
    The scalar $c_N$ is the same for every sector $k$: it is the single global
    proportionality scalar of \cite[Theorem~\texttt{thm1}, lines~1167--1170]%
    {Cirac2016MPDO_arXiv} at length $N$.
\end{lemma}

\begin{proof}\leanok
    For each $N$ on the proportionality tail, expand both total MPVs in their
    sector bases and use the global proportionality scalar $c_N$:
    \begin{align}
        \sum_{j\in J(P)} c_N^{(j)}(P)\ket{V^{(N)}(P_j)}
         & =
        c_N\sum_{k\in J(Q)} c_N^{(k)}(Q)\ket{V^{(N)}(Q_k)}.
        \notag
    \end{align}
    The matched phase identities (\ref{eq:ft_phase_relation_hyp}) give
    $\ket{V^{(N)}(Q_k)}=\zeta_k^N\ket{V^{(N)}(P_{\beta(k)})}$; because every
    sector is matched, reindexing the $Q$-sum by $\beta$ rewrites the right-hand
    side entirely in the $P$-basis. The BNT eventual linear independence of the
    $P$-basis and Lemma~\ref{lem:eventual_coeff_eq_of_eventual_li} then extract
    the coefficient identities (\ref{eq:ft_coeff_identity_threaded}) for all
    sufficiently large $N$.
\end{proof}


\section{Direct-sum conjugation}
\label{sec:app_ft_direct_sum_conjugation}

We now prove the general block-matrix identities used by the global gauge in
Section~\ref{sec:ft_equal_global_gauge}. These statements do not choose or
match sectors; they only combine conjugacies that have already been obtained.

\begin{definition}[Invertible matrix associated to a unitary]
    \label{def:unitary_gl}
    \lean{MPSTensor.unitaryGL}
    \leanok
    A unitary matrix $U$ determines an invertible matrix whose inverse is
    $U^\dagger$.
\end{definition}

\begin{lemma}[Unitary block-diagonal gauges]
    \label{lem:unitary_global_gauge_of_blocks}
    \lean{MPSTensor.globalGaugeOfBlocks_unitaryGL_mem}
    \leanok
    \uses{def:unitary_gl, def:global_gauge_of_blocks}
    If every block gauge $U_k$ is unitary, then the flattened block-diagonal
    gauge $U=\bigoplus_k U_k$ is unitary.
\end{lemma}
\begin{proof}\leanok
    The products are computed blockwise:
    \begin{align}
        U U^\dagger
         & = \bigoplus_k U_kU_k^\dagger
        = \bigoplus_k\mathbb 1
        = \mathbb 1.
        \notag
    \end{align}
\end{proof}

\begin{lemma}[Direct-sum conjugation by the flattened gauge]%
    \label{lem:tensor_from_blocks_globalGauge_conj}
    \lean{MPSTensor.toTensorFromBlocks_eq_globalGaugeOfBlocks_conj,
        MPSTensor.gaugeEquiv_toTensorFromBlocks_of_blockConj}
    \leanok
    \uses{def:global_gauge_of_blocks, def:block_diagonal_tensor, def:gauge_equiv}
    If, for every block $k$ and physical index $i$,
    $B_k^i=X_kA_k^iX_k^{-1}$, then the weighted direct sums satisfy
    \begin{align}
        \bigoplus_k \mu_k B_k^i
         & =
        X\left(\bigoplus_k \mu_k A_k^i\right)X^{-1},
        \notag
    \end{align}
    where $X$ is the flattened block-diagonal gauge from
    Definition~\ref{def:global_gauge_of_blocks}. Hence the two assembled
    tensors are gauge equivalent.
\end{lemma}

\begin{proof}\leanok
    Multiplication of block-diagonal matrices is computed blockwise. The scalar
    weight $\mu_k$ commutes with the conjugation on the $k$th block, and the
    reindexing from dependent block coordinates to the flattened bond index
    preserves products and inverses.
\end{proof}

\begin{theorem}[Global gauge from a copy-weight matching]%
    \label{thm:sector_bnt_global_gauge_of_copy_weight_matching}
    \lean{MPSTensor.sector_bnt_global_gauge_of_copy_weight_matching}
    \leanok
    \uses{def:sector_bnt_copy_weight_matching,
        thm:sector_bnt_global_gauge_of_matched_weights}
    Suppose that the sectors of $P$ and $Q$ are matched, with
    bond-dimension equalities, nonzero phases, and block gauges satisfying
    $Q_k^i=\zeta_k X_kP_{\beta(k)}^iX_k^{-1}$.
    If the same phases also give a copy-weight matching, then the flattened
    $Q$-tensor is obtained from the matched-coordinate $P$-tensor by the
    direct-sum gauge $X=\bigoplus_k(\mathbf{1}_{r_k}\otimes X_k)$.
    This is only a reformulation of
    Theorem~\ref{thm:sector_bnt_global_gauge_of_matched_weights}.
\end{theorem}

\begin{proof}\leanok
    The copy-weight matching supplies
    $\nu_{k,q}=\zeta_k^{-1}\mu_{\beta(k),\tau_k(q)}$. Together with
    $Q_k^i=\zeta_kX_kP_{\beta(k)}^iX_k^{-1}$, these are exactly the two
    hypotheses (\ref{eq:ft_matched_weights_hypotheses}) of
    Theorem~\ref{thm:sector_bnt_global_gauge_of_matched_weights}.
\end{proof}


\section{Additional proportional consequences}
\label{sec:app_ft_conditional_proportional}

The following implications are useful when copy weights satisfy extra
relations. They are not conclusions of CPSV16 Theorem~2.10. In particular,
their additional hypotheses fail for the proportional BNT forms in
Remark~\ref{rem:proportional_length_scalar_gap}.

\begin{theorem}[Additional proportional global gauge under copy-weight matching]%
    \label{thm:sector_bnt_proportional_global_gauge_of_copy_weight_matching}
    \lean{MPSTensor.ft_sector_bnt_proportional_global_gauge_of_copy_weight_matching}
    \leanok
    \uses{def:sector_bnt_copy_weight_matching,
        thm:sector_bnt_proportional_sector_match_witnesses,
        thm:sector_bnt_global_gauge_of_copy_weight_matching}
    Let $P$ and $Q$ be BNT canonical forms whose total tensors generate
    eventually nonzero proportional MPV families.
    Then the proportional sector-matching theorem gives $\beta$, the
    bond-dimension equalities, the unit phases $\zeta_k$, and the block gauges
    $X_k$. If one additionally assumes a copy-weight matching for these same phases, then the flattened
    $Q$-tensor is obtained from the matched-coordinate $P$-tensor by the
    direct-sum gauge $X=\bigoplus_k(\mathbf{1}_{r_k}\otimes X_k)$.
    % Per-sector unit-modulus restriction eliminated by the exact matcher;
    % closure recorded in
    % docs/paper-gaps/cpsv16_global_vs_persector_unit_witness.tex.
\end{theorem}

\begin{proof}\leanok
    The proportional sector-matching theorem supplies $\beta$, $X_k$, and
    unit phases $\zeta_k$ satisfying
    $Q_k^i=\zeta_kX_kP_{\beta(k)}^iX_k^{-1}$. Since $|\zeta_k|=1$, each
    $\zeta_k$ is nonzero. A copy-weight matching gives
    $\nu_{k,q}=\zeta_k^{-1}\mu_{\beta(k),\tau_k(q)}$, so
    Theorem~\ref{thm:sector_bnt_global_gauge_of_copy_weight_matching} gives the
    direct-sum gauge.
\end{proof}

\begin{theorem}[Additional proportional global gauge under coefficient identities]%
    \label{thm:sector_bnt_proportional_global_gauge_of_coeff_identity}
    \lean{MPSTensor.ft_sector_bnt_proportional_global_gauge_of_coeff_identity}
    \leanok
    \uses{lem:sector_bnt_copy_weight_matching_of_coeff_identity,
        thm:sector_bnt_proportional_global_gauge_of_copy_weight_matching}
    Let $P$ and $Q$ be BNT canonical forms whose total tensors generate
    eventually nonzero proportional MPV families.
    Then there are a bijection $\beta:J(Q)\simeq J(P)$, bond-dimension
    equalities, unit-modulus phases $\zeta_k$, and block gauges $X_k$ such that
    $Q_k^i=\zeta_k X_kP_{\beta(k)}^iX_k^{-1}$ for every sector $k$ and physical
    index $i$. If one additionally assumes that these same phases satisfy the matched-sector coefficient
    identities $c_N^{(\beta(k))}(P)=\zeta_k^N\,c_N^{(k)}(Q)$ for all
    sufficiently large $N$, then they determine a copy-weight matching and
    hence the flattened $Q$-tensor is obtained from the matched-coordinate
    $P$-tensor by the direct-sum gauge
    $X=\bigoplus_k(\mathbf{1}_{r_k}\otimes X_k)$.
    % Per-sector unit-modulus restriction eliminated by the exact matcher;
    % closure recorded in
    % docs/paper-gaps/cpsv16_global_vs_persector_unit_witness.tex.
\end{theorem}

\begin{proof}\leanok
    The sector-matching theorem gives $\beta$, $X_k$, and $|\zeta_k|=1$ with
    $Q_k^i=\zeta_kX_kP_{\beta(k)}^iX_k^{-1}$. Since $\zeta_k\ne0$, the
    coefficient identities
    $c_N^{(\beta(k))}(P)=\zeta_k^Nc_N^{(k)}(Q)$ produce
    $\nu_{k,q}=\zeta_k^{-1}\mu_{\beta(k),\tau_k(q)}$ by
    Lemma~\ref{lem:sector_bnt_copy_weight_matching_of_coeff_identity}.
    Theorem~\ref{thm:sector_bnt_proportional_global_gauge_of_copy_weight_matching}
    then applies to this copy-weight matching.
\end{proof}
